\documentclass[10pt]{article}

% Deliberately no font package (e.g. mathptmx, times, fontspec) --
% the lab manual PDF is set in plain Computer Modern at 10pt (verified
% via pdffonts: CMR10/CMBX10/CMBX12/CMMI10/CMSY10/CMEX10/CMTI10), which
% is LaTeX's default. Loading a font package here would deviate from that.

\usepackage[margin=0.8in]{geometry}
\usepackage{graphicx}
\usepackage{amsmath}
\usepackage{booktabs}
\usepackage{float}
\usepackage{hyperref}
\hypersetup{colorlinks=true, linkcolor=black, citecolor=black, urlcolor=blue}

\begin{document}

\begin{center}
{\Large\bfseries Shiv Nadar University Chennai}\\[2pt]
{\bfseries CS3807 -- Deep Learning Laboratory}\\[2pt]
{\bfseries Experiment 3}\\[4pt]
{\large\bfseries Implementation of Convolutional Neural Networks (CNNs) for Image Classification}\\[8pt]
Ankith U Davey \quad 24011101004
\end{center}

\vspace{6pt}
\begin{center}
\begin{tabular}{|p{4cm}|p{9.3cm}|p{2.3cm}|}
\hline
Degree \& Branch & B.Tech Artificial Intelligence \& Data Science & Semester V \\
\hline
Subject Code \& Name & CS3807 -- Deep Learning Laboratory & AY: 2026--27 \\
\hline
\end{tabular}
\end{center}

\vspace{10pt}

\section{Objective}
To understand the working principle of Convolutional Neural Networks by implementing
convolution, pooling, feature map visualization, and image classification using
TensorFlow/Keras.

\section{Background Theory}
A Convolutional Neural Network consists of:
\begin{itemize}
    \item Convolution Layer
    \item Activation Layer
    \item Pooling Layer
    \item Fully Connected Layer
\end{itemize}

The convolution operation is
$$Y(i,j) = \sum_m \sum_n X(i+m, j+n)\,K(m,n)$$
where
\begin{itemize}
    \item $X$ : Input Image
    \item $K$ : Kernel (Filter)
    \item $Y$ : Feature Map
\end{itemize}

The output feature map size is
$$\frac{N - F + 2P}{S} + 1$$
where
\begin{tabular}{@{}ll}
$N$ & Input Size \\
$F$ & Filter Size \\
$P$ & Padding \\
$S$ & Stride \\
\end{tabular}

\subsection{Common Convolution Kernels}
A kernel (or filter) is a small matrix that slides over an input image to extract
important visual features such as edges, textures, corners, and smooth regions. During
convolution, the kernel is multiplied element-wise with the corresponding image
pixels, and the results are summed to produce a single value in the output feature
map. Different kernels highlight different image characteristics.

\subsubsection*{1. Identity Kernel}
The identity kernel preserves the original image without modifying its pixel values.
$$\begin{bmatrix} 0 & 0 & 0 \\ 0 & 1 & 0 \\ 0 & 0 & 0 \end{bmatrix}$$
Purpose:
\begin{itemize}
    \item Preserves the original image.
    \item Used for understanding the convolution operation.
\end{itemize}

\subsubsection*{2. Sobel-X Kernel (Vertical Edge Detection)}
The Sobel-X kernel detects vertical edges by measuring intensity changes along the horizontal direction.
$$\begin{bmatrix} -1 & 0 & 1 \\ -2 & 0 & 2 \\ -1 & 0 & 1 \end{bmatrix}$$
Applications:
\begin{itemize}
    \item Object detection
    \item Lane detection
    \item Face recognition
\end{itemize}

\subsubsection*{3. Sobel-Y Kernel (Horizontal Edge Detection)}
The Sobel-Y kernel detects horizontal edges by measuring intensity changes along the vertical direction.
$$\begin{bmatrix} -1 & -2 & -1 \\ 0 & 0 & 0 \\ 1 & 2 & 1 \end{bmatrix}$$
Applications:
\begin{itemize}
    \item Text detection
    \item Boundary extraction
    \item Medical image analysis
\end{itemize}

\subsubsection*{4. Laplacian Kernel}
The Laplacian kernel detects edges in all directions using the second-order derivative of image intensity.
$$\begin{bmatrix} 0 & -1 & 0 \\ -1 & 4 & -1 \\ 0 & -1 & 0 \end{bmatrix}$$
Applications:
\begin{itemize}
    \item Edge enhancement
    \item Satellite image analysis
    \item Medical imaging
\end{itemize}

\subsubsection*{5. Sharpening Kernel}
This kernel enhances fine details by emphasizing high-frequency components.
$$\begin{bmatrix} 0 & -1 & 0 \\ -1 & 5 & -1 \\ 0 & -1 & 0 \end{bmatrix}$$
Applications:
\begin{itemize}
    \item Image enhancement
    \item Optical Character Recognition (OCR)
    \item Face recognition
\end{itemize}

\subsubsection*{6. Box Blur Kernel}
The Box Blur kernel smooths an image by replacing each pixel with the average of its neighboring pixels.
$$\frac{1}{9}\begin{bmatrix} 1 & 1 & 1 \\ 1 & 1 & 1 \\ 1 & 1 & 1 \end{bmatrix}$$
Applications:
\begin{itemize}
    \item Noise removal
    \item Image smoothing
    \item Image preprocessing
\end{itemize}

\subsubsection*{7. Gaussian Blur Kernel}
Gaussian blur smooths the image while preserving the overall structure better than a simple average filter.
$$\frac{1}{16}\begin{bmatrix} 1 & 2 & 1 \\ 2 & 4 & 2 \\ 1 & 2 & 1 \end{bmatrix}$$
Applications:
\begin{itemize}
    \item Noise reduction
    \item Computer vision preprocessing
    \item Feature extraction
\end{itemize}

\subsubsection*{8. Emboss Kernel}
The Emboss kernel creates a three-dimensional appearance by emphasizing directional changes in intensity.
$$\begin{bmatrix} -2 & -1 & 0 \\ -1 & 1 & 1 \\ 0 & 1 & 2 \end{bmatrix}$$
Applications:
\begin{itemize}
    \item Texture analysis
    \item Artistic image effects
\end{itemize}

\subsubsection*{9. Outline Detection Kernel}
This kernel highlights object boundaries by emphasizing regions with rapid intensity changes.
$$\begin{bmatrix} -1 & -1 & -1 \\ -1 & 8 & -1 \\ -1 & -1 & -1 \end{bmatrix}$$
Applications:
\begin{itemize}
    \item Image segmentation
    \item Boundary detection
    \item Shape extraction
\end{itemize}

\subsubsection*{10. Motion Blur Kernel}
This kernel simulates motion blur along the diagonal direction.
$$\frac{1}{5}\begin{bmatrix}
1 & 0 & 0 & 0 & 0 \\
0 & 1 & 0 & 0 & 0 \\
0 & 0 & 1 & 0 & 0 \\
0 & 0 & 0 & 1 & 0 \\
0 & 0 & 0 & 0 & 1
\end{bmatrix}$$
Applications:
\begin{itemize}
    \item Motion simulation
    \item Image restoration
    \item Video processing
\end{itemize}

\subsubsection*{Summary of Common Kernels}
\begin{table}[H]
\centering
\begin{tabular}{lllll}
\toprule
Kernel & Size & Purpose & Applications \\
\midrule
Identity & $3\times3$ & Preserve original image & Baseline comparison \\
Sobel-X & $3\times3$ & Detect vertical edges & Object detection \\
Sobel-Y & $3\times3$ & Detect horizontal edges & Boundary detection \\
Laplacian & $3\times3$ & Detect edges in all directions & Medical imaging \\
Sharpen & $3\times3$ & Enhance details & Image enhancement \\
Box Blur & $3\times3$ & Smooth image & Noise removal \\
Gaussian Blur & $3\times3$ & Reduce noise & Computer vision \\
Emboss & $3\times3$ & 3D effect & Artistic filtering \\
Outline & $3\times3$ & Boundary extraction & Segmentation \\
Motion Blur & $5\times5$ & Simulate motion & Video processing \\
\bottomrule
\end{tabular}
\caption{Summary of common convolution kernels.}
\end{table}

\subsection{Numerical Example 1: Convolution}
Consider the input image
$$X = \begin{bmatrix} 1 & 2 & 3 \\ 4 & 5 & 6 \\ 7 & 8 & 9 \end{bmatrix}, \quad
K = \begin{bmatrix} 1 & 0 \\ 0 & 1 \end{bmatrix}$$

\noindent Output:
\begin{align*}
\text{First position: } & 1(1) + 2(0) + 4(0) + 5(1) = 6 \\
\text{Second position: } & 2(1) + 3(0) + 5(0) + 6(1) = 8 \\
\text{Third position: } & 4(1) + 5(0) + 7(0) + 8(1) = 12 \\
\text{Fourth position: } & 5(1) + 6(0) + 8(0) + 9(1) = 14
\end{align*}

\noindent Feature Map:
$$\begin{bmatrix} 6 & 8 \\ 12 & 14 \end{bmatrix}$$

\subsection{Numerical Example 2: Max Pooling}
Feature map:
$$\begin{bmatrix} 1 & 5 & 2 & 3 \\ 7 & 8 & 1 & 0 \\ 4 & 6 & 9 & 5 \\ 2 & 3 & 1 & 8 \end{bmatrix}$$

\noindent Using a $2\times2$ max pooling filter:
\begin{align*}
\text{Window 1: } \begin{bmatrix} 1 & 5 \\ 7 & 8 \end{bmatrix} &\rightarrow 8 &
\text{Window 2: } \begin{bmatrix} 2 & 3 \\ 1 & 0 \end{bmatrix} &\rightarrow 3 \\
\text{Window 3: } \begin{bmatrix} 4 & 6 \\ 2 & 3 \end{bmatrix} &\rightarrow 6 &
\text{Window 4: } \begin{bmatrix} 9 & 5 \\ 1 & 8 \end{bmatrix} &\rightarrow 9
\end{align*}

\noindent Output:
$$\begin{bmatrix} 8 & 3 \\ 6 & 9 \end{bmatrix}$$

\subsection{Numerical Example 3: Parameter Calculation}
Suppose

\vspace{4pt}
\noindent Input Image
$$32 \times 32 \times 3$$

\noindent Convolution Layer
\begin{itemize}
    \item Filters $= 16$
    \item Kernel $= 3 \times 3$
\end{itemize}

\noindent Parameters
$$(3 \times 3 \times 3 + 1) \times 16$$
$$= (27 + 1) \times 16$$
$$= 448$$

\noindent Therefore, total trainable parameters $= 448$.

\section{Dataset}
Dataset: CIFAR-10
\begin{itemize}
    \item Training Images: $50{,}000$
    \item Testing Images: $10{,}000$
    \item Classes: $10$
    \item Image Size: $32 \times 32 \times 3$
\end{itemize}

\section{Source Code}
The complete source code is available in the accompanying GitHub repository:
\url{https://github.com/Ankith-Davey/Deep-Learning-Lab-Ankith-24011101004/tree/main/Lab%203}

\section{Task 1: Load CIFAR-10, Display Samples, Print Dimensions, Plot Class Distribution}
Load CIFAR-10.
\begin{itemize}
    \item Display ten sample images.
    \item Print dataset dimensions.
    \item Plot class distribution.
\end{itemize}

The dataset was loaded and verified directly:
\begin{itemize}
    \item \texttt{x\_train} shape: $(50000, 32, 32, 3)$
    \item \texttt{y\_train} shape: $(50000,)$
    \item \texttt{x\_test} shape: $(10000, 32, 32, 3)$
    \item \texttt{y\_test} shape: $(10000,)$
\end{itemize}

\begin{figure}[H]
    \centering
    \includegraphics[width=0.85\textwidth]{01_sample_images.png}
    \caption{Ten randomly sampled CIFAR-10 training images with their class labels.}
    \label{fig:sample_images}
\end{figure}

\textbf{Inference:} The ten sampled images span natural scenes (frog, deer), animals
(cat, dog), and vehicles (truck, ship, automobile) at low $32\times32$ resolution,
confirming the dataset loaded correctly with intact labels. Even at this resolution,
class-defining shapes remain visually distinguishable to a human eye.

\begin{figure}[H]
    \centering
    \includegraphics[width=0.7\textwidth]{02_class_distribution.png}
    \caption{Class distribution across the training set.}
    \label{fig:class_dist}
\end{figure}

\textbf{Inference:} All ten classes are represented by exactly $5{,}000$ training
images, confirming CIFAR-10 is perfectly class-balanced. This matters directly for
evaluation: overall accuracy is not misleading the way it would be for an imbalanced
dataset, so it can be interpreted at face value.

\section{Task 2: Convolution -- Kernel Size Comparison}
Implement a convolution layer.

\noindent Compare
\begin{itemize}
    \item Kernel $= 3 \times 3$
    \item Kernel $= 5 \times 5$
    \item Kernel $= 7 \times 7$
\end{itemize}

\noindent Record feature map sizes.

\vspace{4pt}
\noindent The output feature map size follows $\left\lfloor \dfrac{N - F + 2P}{S} \right\rfloor + 1$,
where $N$ is the input size, $F$ the filter size, $P$ the padding, and $S$ the stride.
Applying $3\times3$, $5\times5$, and $7\times7$ kernels to a $32\times32$ input with
valid padding and stride $1$ gave the following, confirmed against actual
\texttt{Conv2D} layer outputs:

\begin{table}[H]
\centering
\begin{tabular}{lllll}
\toprule
Kernel & Padding & Stride & Output Size & Feature Map Shape \\
\midrule
$3\times3$ & Valid (0) & 1 & $30\times30$ & $(30, 30, 8)$ \\
$5\times5$ & Valid (0) & 1 & $28\times28$ & $(28, 28, 8)$ \\
$7\times7$ & Valid (0) & 1 & $26\times26$ & $(26, 26, 8)$ \\
\bottomrule
\end{tabular}
\caption{Feature map size vs. kernel size (32$\times$32 input, valid padding, stride 1).}
\end{table}

\textbf{Inference:} Larger kernels shrink the output feature map faster under valid
padding, since more border pixels are lost per convolution pass. This is a direct
trade-off: a $7\times7$ kernel covers more spatial context per step but discards more
information at the edges and produces smaller downstream feature maps than a $3\times3$
kernel would after the same number of layers.

\section{Task 3: Hyperparameter Study -- Stride and Padding}
Study hyperparameters.

\noindent Compare
\begin{itemize}
    \item Stride $= 1$
    \item Stride $= 2$
    \item Padding $=$ Same
    \item Padding $=$ Valid
\end{itemize}

\noindent Compute output dimensions.

\begin{table}[H]
\centering
\begin{tabular}{lll}
\toprule
Stride & Padding & Output Shape ($F=3$) \\
\midrule
1 & Valid & $(30, 30, 8)$ \\
1 & Same  & $(32, 32, 8)$ \\
2 & Valid & $(15, 15, 8)$ \\
2 & Same  & $(16, 16, 8)$ \\
\bottomrule
\end{tabular}
\caption{Output shape as a function of stride and padding, kernel fixed at $3\times3$.}
\end{table}

\textbf{Inference:} ``Same'' padding keeps the output spatial size equal to (stride 1)
or a clean half (stride 2) of the input size, since it pads the input specifically to
compensate for the pixels a valid convolution would otherwise lose. Stride has a much
larger effect than padding on output size: doubling the stride roughly halves each
spatial dimension regardless of the padding choice.

\section{Task 4: Feature Map Visualization}
Feature maps from the first convolution layer are shown here both before and after
training, for direct comparison, using the same \texttt{conv1} layer and the same
input image in each case.

\begin{figure}[H]
    \centering
    \includegraphics[width=0.85\textwidth]{07a_feature_maps_untrained.png}
    \caption{Feature maps from the first convolution layer, untrained (random) weights.}
    \label{fig:fmap_untrained}
\end{figure}

\textbf{Inference:} Even with random weights, most filters still trace the input
image's silhouette and background texture, since a random linear mix of RGB channels
partially preserves edge structure. The maps are not pure noise, but none show the
sharp, filter-specific selectivity that trained weights produce.

\begin{figure}[H]
    \centering
    \includegraphics[width=0.85\textwidth]{07b_feature_maps_trained.png}
    \caption{Feature maps from the first convolution layer, after training the CNN classifier.}
    \label{fig:fmap_trained}
\end{figure}

\textbf{Inference:} Several filters now show sharper, thinner responses to specific
edges and contours compared to the blobbier untrained versions (e.g. Filters 3, 6, and
7). This indicates the first layer learned genuine low-level edge detectors rather than
passing through unrefined input structure.

\section{Task 5: Max Pooling vs. Average Pooling}
Compare
\begin{itemize}
    \item Max Pooling
    \item Average Pooling
\end{itemize}

\noindent Compare output size and accuracy.

\vspace{4pt}
\noindent Pooling output size and validation accuracy were compared separately, as follows:

\begin{table}[H]
\centering
\begin{tabular}{lll}
\toprule
Layer & Max Pooling Output Shape & Average Pooling Output Shape \\
\midrule
Pool 1 & $(16, 16, 16)$ & $(16, 16, 16)$ \\
Pool 2 & $(8, 8, 32)$   & $(8, 8, 32)$ \\
\bottomrule
\end{tabular}
\caption{Output shape comparison: max vs. average pooling, identical kernel/stride.}
\end{table}

\begin{table}[H]
\centering
\begin{tabular}{lll}
\toprule
Pooling Type & Final Validation Accuracy (10 epochs) & Training Time \\
\midrule
Max Pooling     & 66.42\% & 61.2\,s \\
Average Pooling & 65.98\% & 59.4\,s \\
\bottomrule
\end{tabular}
\caption{Max vs. average pooling: accuracy and training time.}
\end{table}

\begin{figure}[H]
    \centering
    \includegraphics[width=0.7\textwidth]{09_pooling_comparison.png}
    \caption{Validation accuracy per epoch: max pooling vs. average pooling.}
    \label{fig:pooling_comparison}
\end{figure}

\textbf{Inference:} Max pooling tracks slightly above average pooling in validation
accuracy across all 10 epochs (66.42\% vs. 65.98\% final), though the gap stays narrow
throughout. Both curves rise with a similar shape, indicating pooling type has a real
but secondary effect compared to other architecture choices.

\section{Task 6: CNN Construction and Training}

$$\text{Input} \rightarrow \text{Conv} \rightarrow \text{ReLU} \rightarrow \text{MaxPool}
\rightarrow \text{Conv} \rightarrow \text{ReLU} \rightarrow \text{MaxPool} \rightarrow
\text{Flatten} \rightarrow \text{Dense} \rightarrow \text{Softmax}$$

\noindent Train for
\begin{itemize}
    \item Optimizer : Adam
    \item Epochs : 20
    \item Batch Size : 32
\end{itemize}

\noindent Total trainable parameters: $25{,}578$.

\begin{figure}[H]
    \centering
    \includegraphics[width=0.7\textwidth]{03_training_accuracy.png}
    \caption{Training accuracy across 20 epochs.}
\end{figure}

\textbf{Inference:} Accuracy rises steadily and near-monotonically across all 20
epochs, reaching $77.69\%$ by the final epoch. The absence of any plateau suggests the
model was still actively fitting the training set when training stopped.

\begin{figure}[H]
    \centering
    \includegraphics[width=0.7\textwidth]{04_validation_accuracy.png}
    \caption{Validation accuracy across 20 epochs.}
\end{figure}

\textbf{Inference:} Accuracy climbs until roughly epoch 12--14 ($\approx 68\%$), then
flattens and drifts slightly downward to $67.45\%$ by epoch 20. This divergence from
the still-rising training curve is the clearest sign of overfitting in the run.

\begin{figure}[H]
    \centering
    \includegraphics[width=0.7\textwidth]{05_training_loss.png}
    \caption{Training loss across 20 epochs.}
\end{figure}

\textbf{Inference:} Loss decreases smoothly and continuously through all 20 epochs
with no sign of instability. This confirms the Adam optimizer converged cleanly on the
training set at this learning rate and batch size.

\begin{figure}[H]
    \centering
    \includegraphics[width=0.7\textwidth]{06_validation_loss.png}
    \caption{Validation loss across 20 epochs.}
\end{figure}

\textbf{Inference:} Loss reaches its minimum around epoch 10--12, then rises for the
remainder of training even as training loss keeps falling. This loss divergence is a
more sensitive overfitting signal than the accuracy curves alone.

\section{Task 7: Evaluation}
Evaluate
\begin{itemize}
    \item Accuracy
    \item Precision
    \item Recall
    \item F1-score
    \item Confusion Matrix
    \item Classification Report
\end{itemize}

\begin{table}[H]
\centering
\begin{tabular}{ll}
\toprule
Metric & Value \\
\midrule
Test Accuracy      & 67.45\% \\
Precision (macro)  & 68.23\% \\
Recall (macro)     & 67.45\% \\
F1-score (macro)   & 67.40\% \\
\bottomrule
\end{tabular}
\caption{Overall test-set evaluation metrics.}
\end{table}

\begin{table}[H]
\centering
\begin{tabular}{lccc}
\toprule
Class & Precision & Recall & F1-score \\
\midrule
airplane   & 0.75 & 0.64 & 0.69 \\
automobile & 0.80 & 0.82 & 0.81 \\
bird       & 0.52 & 0.62 & 0.57 \\
cat        & 0.55 & 0.38 & 0.45 \\
deer       & 0.54 & 0.69 & 0.61 \\
dog        & 0.58 & 0.62 & 0.60 \\
frog       & 0.72 & 0.75 & 0.73 \\
horse      & 0.69 & 0.77 & 0.73 \\
ship       & 0.83 & 0.75 & 0.79 \\
truck      & 0.84 & 0.70 & 0.76 \\
\bottomrule
\end{tabular}
\caption{Per-class precision, recall, and F1-score.}
\end{table}

\begin{figure}[H]
    \centering
    \includegraphics[width=0.75\textwidth]{08_confusion_matrix.png}
    \caption{Confusion matrix on the CIFAR-10 test set.}
    \label{fig:confusion}
\end{figure}

\textbf{Inference:} Cat is the weakest class by far -- only 379 correctly classified
out of 1,000, with 213 confused as dog and 136 as deer. Vehicle classes show far
cleaner diagonals (automobile: 822, ship: 750), confirming errors concentrate among
visually similar animal classes rather than being spread evenly.

\section{Additional Exercises}

\subsection*{Exercise 1: Output Size Calculation}
For a $64\times64$ input, $5\times5$ kernel, stride $2$, padding $2$:
$$\left\lfloor \frac{64 - 5 + 2(2)}{2} \right\rfloor + 1 = \left\lfloor \frac{63}{2} \right\rfloor + 1 = 31 + 1 = 32$$
\textbf{Output size: $32 \times 32$.}

\subsection*{Exercise 2: Trainable Parameters}
For a convolution layer with $64$ filters of size $3\times3$ on an RGB (3-channel) input:
$$(3 \times 3 \times 3 + 1) \times 64 = 28 \times 64 = 1{,}792 \text{ parameters}$$

\subsection*{Exercise 3: ReLU vs. Sigmoid Activation}

\begin{table}[H]
\centering
\begin{tabular}{ll}
\toprule
Activation & Final Validation Accuracy (5 epochs) \\
\midrule
ReLU    & 63.81\% \\
Sigmoid & 47.83\% \\
\bottomrule
\end{tabular}
\caption{ReLU vs. Sigmoid activation, same architecture, 5 epochs each.}
\end{table}

\begin{figure}[H]
    \centering
    \includegraphics[width=0.7\textwidth]{10_relu_vs_sigmoid.png}
    \caption{Validation accuracy per epoch: ReLU vs. Sigmoid.}
\end{figure}

\textbf{Inference:} ReLU pulls ahead within the first epoch and maintains a wide,
growing lead throughout, ending at $63.81\%$ vs. Sigmoid's $47.83\%$. This matches
Sigmoid's known vanishing-gradient weakness, which slows early-layer learning under a
short 5-epoch budget.

\subsection*{Exercise 4: Max Pooling vs. Average Pooling}
Max pooling and average pooling were compared using the same architecture, trained
for 10 epochs each. Max pooling reached $66.42\%$ final validation accuracy versus
$65.98\%$ for average pooling. Both produced identical output shapes at every pooling
layer, since output size under pooling depends only on kernel size and stride, not on
the aggregation function used.

\subsection*{Exercise 5: Filter Count -- 16 vs. 64}

\begin{table}[H]
\centering
\begin{tabular}{lll}
\toprule
Filter Configuration & Validation Accuracy (5 epochs) & Training Time \\
\midrule
16 / 32 filters & 64.59\% & 39.2\,s \\
64 / 128 filters & 68.04\% & 41.8\,s \\
\bottomrule
\end{tabular}
\caption{Effect of filter count on accuracy and computation time.}
\end{table}

\textbf{Inference:} Increasing filter counts from 16/32 to 64/128 improved validation
accuracy by roughly 3.5 points within the same 5-epoch budget, at a modest training-time
cost (about 2.6 seconds more per full run, not per epoch). More filters let the network
learn a richer, more diverse set of low- and mid-level feature detectors in the same
number of layers, which explains the accuracy gain; the relatively small time increase
reflects that convolution cost scales with filter count roughly linearly, while
overall training time on a GPU is also bounded by data-loading and other fixed
overheads that don't scale with filter count.

\section{Results}
Record
\begin{itemize}
    \item Training Accuracy
    \item Testing Accuracy
    \item Precision
    \item Recall
    \item F1-score
    \item Number of Parameters
\end{itemize}

\begin{table}[H]
\centering
\begin{tabular}{ll}
\toprule
Metric & Value \\
\midrule
Training Accuracy (final epoch) & 77.69\% \\
Testing Accuracy                & 67.45\% \\
Precision (macro)               & 68.23\% \\
Recall (macro)                  & 67.45\% \\
F1-score (macro)                & 67.40\% \\
Total Trainable Parameters      & 25{,}578 \\
\bottomrule
\end{tabular}
\caption{Summary of final results, CNN classifier.}
\end{table}

\section{Discussion}

\subsection*{1. Why is convolution preferred over fully connected layers for images?}
Convolution layers use a small, spatially-shared kernel that slides across the entire
image, so the same set of weights detects a given feature (an edge, a texture)
regardless of where it appears in the image. A fully connected layer, by contrast,
assigns an independent weight to every input pixel, which both discards this spatial
locality/translation-invariance structure and produces vastly more parameters for the
same input size -- for a $32\times32\times3$ image, a single fully connected layer to
even a modest hidden size would need orders of magnitude more parameters than the
$25{,}578$ used by the entire CNN classifier built in this experiment.

\subsection*{2. How does stride affect the feature map size?}
Stride controls how far the kernel moves between applications. The stride and padding
comparison confirmed this directly: a stride of $2$ roughly halves each spatial
dimension of the output compared to a stride of $1$, since the kernel visits half as
many positions. Larger strides therefore produce smaller feature maps and reduce
computation, at the cost of skipping some spatial positions the kernel would otherwise
have covered.

\subsection*{3. What is the role of padding?}
Padding adds extra border pixels (typically zeros) around the input before
convolution. Its main role, confirmed by the same stride and padding comparison, is to
control output size: ``same'' padding is chosen specifically so the output retains the
input's spatial dimensions (at stride 1), which allows deeper stacks of convolution
layers without the feature map shrinking away entirely. It also lets border pixels
participate in more convolution windows than they would under valid padding, where
they contribute to comparatively fewer output values.

\subsection*{4. Why is pooling used?}
Pooling reduces the spatial dimensions of feature maps (the max pooling vs. average
pooling comparison showed both pooling
types roughly halved each spatial dimension), which reduces the number of parameters
and computation needed in subsequent layers and provides a degree of translation
invariance -- small shifts in where a feature appears in the input have less effect on
the pooled output. It also acts as a mild form of regularization by discarding precise
spatial information that isn't necessary for classification.

\subsection*{5. How do feature maps represent image characteristics?}
Each filter in a convolution layer learns to respond strongly to a specific visual
pattern -- an edge orientation, a colour transition, a texture. The feature map
visualization confirmed this directly: an untrained filter's feature map is close to
random noise, while a trained filter's feature map highlights the specific regions of
the input image where that learned pattern is present. Stacking convolution layers
builds increasingly abstract representations: early layers (like the \texttt{conv1}
layer visualized here) tend to capture low-level features such as edges and colour
boundaries, while deeper layers combine these into higher-level, more class-specific
patterns.

\subsection*{6. Why do CNNs require fewer parameters than MLPs?}
CNNs exploit two structural properties of images that MLPs ignore: local connectivity
(each output depends only on a small neighborhood of the input, not the entire image)
and weight sharing (the same kernel is reused across every spatial position). As
computed in Exercise 2, a $3\times3$ convolution layer with 64 filters over an RGB
input needs only $1{,}792$ parameters, regardless of the input's spatial size -- an MLP
layer connecting even a modest number of input pixels to a hidden layer of comparable
width would need vastly more, since every input-output pair gets its own independent
weight.

\section{Conclusion}
This experiment implemented and evaluated the core building blocks of a Convolutional
Neural Network -- convolution, pooling, and a full classification pipeline -- on the
CIFAR-10 dataset, using the specified architecture with no additional layers. The
resulting model, with $25{,}578$ trainable parameters,
reached $77.69\%$ training accuracy and $67.45\%$ test accuracy after 20 epochs. The
gap between these two numbers, together with the training/validation loss curves for
the main CNN classifier, points clearly to overfitting: the architecture includes no
regularization (no dropout, batch normalization, or data augmentation), so the model's
capacity, while small, is spent partly on memorizing training-set specifics rather than
generalizable features, particularly after epoch 12--14. The hyperparameter
comparisons reinforced expected theoretical behaviour throughout: max and average
pooling produced identical output shapes while differing slightly in accuracy, ReLU
substantially outperformed Sigmoid under a limited epoch budget, and increasing filter
count improved accuracy at a modest computational cost. Per-class evaluation further
showed that visually distinctive classes (automobile, ship, truck) were classified far
more reliably than visually similar ones (cat, bird, dog), indicating that the model's
errors are concentrated in genuinely difficult, closely related categories rather than
being spread uniformly across all classes. Overall, the experiment met its stated
objective of building working intuition for how convolution, activation, and pooling
operations combine into a trainable image classifier, while also surfacing a concrete,
observable example of overfitting under a deliberately unregularized architecture.

\section{References}
\begin{enumerate}
    \item Goodfellow, I., Bengio, Y., \& Courville, A. \textit{Deep Learning}. MIT Press.
    \item Bishop, C. M. \textit{Pattern Recognition and Machine Learning}. Springer.
    \item Haykin, S. \textit{Neural Networks and Learning Machines}. Pearson.
    \item TensorFlow Documentation. \url{https://www.tensorflow.org/}
    \item CIFAR-10 Dataset Documentation. \url{https://www.cs.toronto.edu/~kriz/cifar.html}
\end{enumerate}

\end{document}